% Options for packages loaded elsewhere
\PassOptionsToPackage{unicode}{hyperref}
\PassOptionsToPackage{hyphens}{url}
\documentclass[
]{article}
\usepackage{xcolor}
\usepackage{amsmath,amssymb}
% \setcounter{secnumdepth}{-\maxdimen} % remove section numbering
\usepackage{iftex}
\ifPDFTeX
  \usepackage[T1]{fontenc}
  \usepackage[utf8]{inputenc}
  \usepackage{textcomp} % provide euro and other symbols
\else % if luatex or xetex
  \usepackage{unicode-math} % this also loads fontspec
  \defaultfontfeatures{Scale=MatchLowercase}
  \defaultfontfeatures[\rmfamily]{Ligatures=TeX,Scale=1}
\fi
\usepackage{lmodern}
\ifPDFTeX\else
  % xetex/luatex font selection
\fi
% Use upquote if available, for straight quotes in verbatim environments
\IfFileExists{upquote.sty}{\usepackage{upquote}}{}
\IfFileExists{microtype.sty}{% use microtype if available
  \usepackage[]{microtype}
  \UseMicrotypeSet[protrusion]{basicmath} % disable protrusion for tt fonts
}{}
\makeatletter
\@ifundefined{KOMAClassName}{% if non-KOMA class
  \IfFileExists{parskip.sty}{%
    \usepackage{parskip}
  }{% else
    \setlength{\parindent}{0pt}
    \setlength{\parskip}{6pt plus 2pt minus 1pt}}
}{% if KOMA class
  \KOMAoptions{parskip=half}}
\makeatother
\usepackage{longtable,booktabs,array}
\newcounter{none} % for unnumbered tables
\usepackage{calc} % for calculating minipage widths
% Correct order of tables after \paragraph or \subparagraph
\usepackage{etoolbox}
\makeatletter
\patchcmd\longtable{\par}{\if@noskipsec\mbox{}\fi\par}{}{}
\makeatother
% Allow footnotes in longtable head/foot
\IfFileExists{footnotehyper.sty}{\usepackage{footnotehyper}}{\usepackage{footnote}}
\makesavenoteenv{longtable}
\setlength{\emergencystretch}{3em} % prevent overfull lines
\providecommand{\tightlist}{%
  \setlength{\itemsep}{0pt}\setlength{\parskip}{0pt}}
\usepackage{bookmark}
\IfFileExists{xurl.sty}{\usepackage{xurl}}{} % add URL line breaks if available
\urlstyle{same}
\hypersetup{
  hidelinks,
  pdfcreator={LaTeX via pandoc}}

\title{From Deltas to Decisions: DeltaRAG and Ecosystem Change\\Intelligence for Digital Risk Monitoring}
\author{}
\date{}

\begin{document}

\maketitle

\begin{abstract}
Operational risk controls in digital ecosystems increasingly depend on
external, fast-changing signals, security advisories, vulnerability
catalogs, API release notes, and platform policy updates, to calibrate
fraud detection thresholds, enforce verification gates, and escalate
incidents. However, monitoring these sources remains predominantly
manual, reactive, and disconnected from downstream decision workflows.
When detection lags, enforcement logic drifts: teams either over-block
legitimate activity or under-block abuse, widening the exposure window
between an ecosystem change and an operational response. We introduce
\textbf{Ecosystem Change Intelligence (ECI)}, a framework that reframes
ecosystem monitoring as a change-centric decision pipeline. ECI treats
each monitored source as a versioned stream, detects deltas between
snapshots, and converts them into structured Change Events that flow
through AI-driven triage and coordination agents to produce
evidence-backed Action Tickets. Methodologically, we contribute
\textbf{DeltaRAG}, a retrieval methodology that combines (a)
delta-grounded vector retrieval on change-local context windows rather
than full documents, and (b) a knowledge graph (KG) augmentation layer
that links change events to entities, vulnerabilities, and cross-source
signals so that escalation decisions can be justified using connected
evidence rather than isolated snippets, the two components fused via
Reciprocal Rank Fusion. We instantiate ECI on the Android digital risk
ecosystem, monitoring security bulletins, Play Integrity API updates,
CISA Known Exploited Vulnerabilities, and developer policy changes.
A 100-query (85 scored) deterministic ablation study across
768-dimensional Nomic v1.5 embeddings stored in Supabase/pgvector
shows that DeltaRAG outperforms Standard RAG on top-rank quality
(P@1\,=\,0.882 vs.\ 0.847; MRR\,=\,0.930 vs.\ 0.913), is statistically
tied on nDCG@5 ($p = 0.43$), and strictly dominates a KG-only baseline
on every metric ($\Delta$R@5\,=\,$+0.328$, $p < 0.001$, Holm-corrected),
while rejecting 100\% of out-of-scope queries. The system achieves
reliable structured-output generation from lightweight LLMs (Llama 3.1
8B) within a multi-agent architecture (Sentinel triage + Coordinator
synthesis) that surfaces prioritized, traceable recommendations to
fraud operations teams.
\end{abstract}


\noindent\textbf{Keywords:} Retrieval-Augmented Generation, Change Intelligence,
Digital Risk Monitoring, Multi-Agent Systems, Knowledge Graphs, Android
Security, Fraud Operations

\section{Introduction}\label{introduction}

Modern digital risk operations depend on a continuous flow of ecosystem
intelligence, security patches, platform policy updates, API
deprecations, and vulnerability disclosures to decide when to adjust
fraud detection models, step up verification, tighten enforcement gates,
or escalate incidents. The Android ecosystem alone publishes monthly
security bulletins, quarterly Play Integrity API revisions, rolling
developer policy updates, and references dozens of CISA Known Exploited
Vulnerabilities (KEV) entries annually. Each of these signals can
silently alter the trust assumptions underlying operational fraud
controls.

Despite the operational criticality of these signals, current monitoring
approaches are predominantly static and rule-based. Teams rely on
keyword alerts, manual reading of changelogs, or brittle heuristics that
treat documents as fixed reference material. In practice, ecosystem
updates are heterogeneous in format, cadence, and terminology, and
``meaningful change'' is rarely captured by simple text diffs. More
importantly, static systems fail at the step that matters most for
operations: \emph{translating updates into auditable action}. Without
evidence-backed outputs that specify what to do next, who should own it,
how urgent it is, and exactly what evidence supports the recommendation,
decision-making remains subjective and difficult to operationalize.

Retrieval-Augmented Generation (RAG) has emerged as a promising paradigm
for grounding large language model (LLM) outputs in retrieved evidence,
reducing hallucination and improving factual accuracy. However, standard
RAG pipelines are designed for static document collections and struggle
with the change-centric nature of ecosystem monitoring. When a security
bulletin is revised or a policy page is updated, standard RAG either
retrieves the entire document (diluting the signal) or misses the change
entirely because the embedding space does not distinguish between old
and new content.

We introduce \textbf{Ecosystem Change Intelligence (ECI)}, a framework
that reframes ecosystem monitoring as a structured decision pipeline.
ECI treats each monitored source as a versioned stream and processes them through an six-stage pipeline implementation
of these three stages: (1) \textbf{Delta Detection}, identifying
meaningful changes between time-stamped snapshots; (2)
\textbf{Change-Grounded Retrieval}, embedding and retrieving
change-local context rather than full documents; and (3)
\textbf{Evidence-Backed Action Generation}, producing structured
recommendations with traceable citations, priority levels, and ownership
assignments.

To implement ECI, we contribute \textbf{DeltaRAG}, a retrieval
methodology that combines two complementary components fused via
Reciprocal Rank Fusion. The first component is a delta-grounded vector
index: changed content (additions and deletions, rather than full
documents) is chunked, embedded using task-aware models, and stored as
versioned vectors, enabling a Sentinel Triage Agent to score relevance
and risk without hallucinating on stale content. The second component
is a \textbf{knowledge graph (KG) augmentation} layer that links change
events to extracted entities (CVEs, APIs, permissions, policy clauses)
via deterministic regex-based extraction and 2-hop graph traversal,
enabling a Coordinator Agent to synthesize cross-source patterns,
identify cascading risks, and produce recommendations justified by
connected evidence rather than isolated snippets. Unlike Graph-RAG
approaches (Edge et al., 2024) that construct entity-relation graphs
from unstructured text using LLMs, our KG layer prioritizes latency,
determinism, and zero hallucination, design choices we revisit in
Section~\ref{discussionn} as both a strength and a future-work target.

We instantiate ECI on the Android digital risk ecosystem, a domain
chosen for its operational significance to fraud teams at financial
institutions and its rich, publicly available signal landscape. Our
system monitors four source categories, Android Security Bulletins,
Play Integrity API documentation, CISA KEV feeds, and Google Developer
Policy updates, and processes them through an six-stage pipeline
from source monitoring to interactive dashboard delivery. The pipeline
uses Nomic v1.5 embeddings (768-dimensional, task-aware, locally
hosted) stored in Supabase with pgvector for sub-10ms similarity
search, and employs lightweight LLMs (Llama 3.1 8B) via Groq for
rapid prototyping with a planned migration to self-hosted GCP
infrastructure.

The primary contribution of this paper is \textbf{DeltaRAG}, the
fused delta-grounded retrieval methodology described above, which
enables high-precision triage of evolving platform signals. Beyond
DeltaRAG, we contribute: (1) the \textbf{ECI framework}, formalizing
ecosystem monitoring as a change-centric decision pipeline with formal
definitions for delta detection, retrieval, and risk scoring; (2) a \textbf{knowledge graph augmentation} design
for cross-source evidence linking via 2-hop traversal with entity-type
weighting and a cross-category bonus; (3) a
\textbf{multi-agent architecture} (Sentinel triage + Coordinator
synthesis) that converts change events into auditable Action Tickets
with formal composite risk scoring (Equation~11); and (4) an
\textbf{ablation study} across 100 deterministic benchmark queries
(85 scored, 10 query-type categories, Holm-corrected significance)
showing that DeltaRAG outperforms Standard RAG on top-rank quality
(P@1\,=\,0.882 vs.\ 0.847; MRR\,=\,0.930 vs.\ 0.913), is statistically
tied on nDCG@5 ($p = 0.43$), and strictly dominates a KG-only baseline
on every metric ($\Delta$R@5\,=\,$+0.328$, $p < 0.001$), while
empirically identifying pure-regex entity extraction as the primary
bottleneck in the standalone KG variant.

\section{Related Work}\label{related-work}

\subsection{Retrieval-Augmented
Generation}\label{retrieval-augmented-generation}

Retrieval-Augmented Generation (RAG) was introduced by Lewis et al.
(2020) as a method for augmenting parametric language models with
non-parametric retrieval over external knowledge stores. The RAG
paradigm has since been widely adopted for knowledge-intensive tasks,
with subsequent work focusing on improving retrieval quality (Izacard \&
Grave, 2021), chunk-level optimization (Shi et al., 2023), and
domain-specific adaptations for legal (Savelka et al., 2023), medical
(Xiong et al., 2024), and cybersecurity (Aghaei et al., 2023)
applications. However, all of these approaches assume a static document
corpus: the same document is retrieved regardless of whether it was
updated an hour ago or a year ago. This static assumption is
fundamentally incompatible with ecosystem monitoring, where the
\emph{change itself},not the document,is the unit of interest.
\textbf{DeltaRAG} closes this gap by restricting the retrieval index to
change-local context (Eq.~5), ensuring that retrieval precision is not
diluted by unchanged background content.

\subsection{Knowledge Graphs for Information
Retrieval}\label{knowledge-graphs-for-information-retrieval}

Knowledge graphs have been integrated with language models in several
retrieval paradigms. Graph-RAG approaches (Edge et al., 2024; He et al.,
2024) construct entity-relation graphs from unstructured text and use
graph traversal to improve retrieval coverage and reasoning depth.
KG-augmented QA systems (Yasunaga et al., 2021) leverage structured
triples to resolve multi-hop queries. In the cybersecurity domain,
knowledge graphs have been used to link CVE entries, attack patterns,
and affected software (Mittal et al., 2016; Li et al., 2022). These
systems operate on static knowledge bases. Our Graph-RAG layer
(Eqs.~6--10) is specifically designed for \emph{change-event linking}:
nodes represent individual delta events, not documents, and edges encode
operational relationships (e.g., a CVE \emph{affects} an API level that
\emph{references} a policy clause). This change-centric graph topology
surfaces cascading risk patterns that would remain invisible to either
document-level or purely vector-based retrieval.

\subsection{AI Agents and Multi-Agent
Systems}\label{ai-agents-and-multi-agent-systems}

The emergence of LLM-powered agents,systems that can plan, use tools,
and take actions,has accelerated with frameworks like ReAct (Yao et
al., 2023), LangChain, and LangGraph. Multi-agent architectures extend
this to cooperative or hierarchical agent teams, with recent work on
AutoGen (Wu et al., 2023) and CAMEL (Li et al., 2023) demonstrating
effective task decomposition. In fraud detection and cybersecurity,
agentic systems have been proposed for threat hunting (Ferrag et al.,
2024) and automated incident response. Prior multi-agent work focuses on
open-ended reasoning and task planning; it does not target
\emph{structured-output stability}, where every agent response must
parse as valid JSON and conform to a fixed operational schema. Our
Sentinel + Coordinator architecture is explicitly optimized for this
requirement, enforcing schema compliance through few-shot examples,
chain-of-thought prefixes, and Pydantic validation,a design choice
validated by our LLM selection rationale in Section~\ref{llm-evaluation-and-selection}.

\subsection{Digital Risk and Threat Intelligence
Monitoring}\label{digital-risk-and-threat-intelligence-monitoring}

Threat intelligence platforms (TIPs) such as MISP, OpenCTI, and
commercial solutions aggregate indicators of compromise (IoCs) and
vulnerability data. However, these systems primarily operate on
structured feeds (STIX/TAXII, CVE JSON) and do not perform semantic
analysis of unstructured platform documentation changes. Platform change
monitoring in industry remains largely manual, with security teams
subscribing to mailing lists and RSS feeds. Recent work on automated
change detection for web content (Patel et al., 2023) and policy
tracking (Li et al., 2024) addresses source-level detection but stops
at the detection boundary: neither system converts detected changes into
risk-scored, evidence-grounded, owner-assigned action items. ECI closes
this last-mile gap, linking detection directly to operational decision
support through DeltaRAG retrieval and the composite risk score
(Eq.~11).

\subsection{Embedding Models for Specialized
Domains}\label{embedding-models-for-specialized-domains}

The choice of embedding model significantly impacts retrieval quality in
RAG systems. General-purpose models like Sentence-BERT (Reimers \&
Gurevych, 2019) and OpenAI embeddings perform well on broad tasks but
may underperform on domain-specific terminology. Task-aware models like
Nomic Embed v1.5 (Nussbaum et al., 2024) allow prefix-based optimization
(\texttt{search\_document:}, \texttt{search\_query:}) that adapts the embedding
space to retrieval use cases. Domain-specialized models like BGE (Xiao
et al., 2024) and cybersecurity-specific embeddings (Aghaei et al.,
2023) trade generality for improved in-domain performance. We benchmark
three model families against a golden dataset of Android security
queries (Section~7.1) and find that \emph{task-awareness} and
\emph{long-context support} (8,192 tokens) are more predictive of
retrieval quality than raw model size,motivating the selection of
Nomic v1.5 as the DeltaRAG backbone.

\section{Problem Formulation: Ecosystem Change
Intelligence}\label{problem-formulation-ecosystem-change-intelligence}

\subsection{The Operational Gap}\label{the-operational-gap}

We define the operational problem as follows. Let
$S = \{s_1, s_2, \ldots, s_n\}$
be a set of monitored ecosystem sources. Each source $s_i$ produces a
time-ordered sequence of snapshots:
$s_i(t_1), s_i(t_2), \ldots$.

A \textbf{change event} $\Delta_i(t)$ is a structured record capturing the
semantic difference between consecutive snapshots
$s_i(t_k)$ and $s_i(t_{k-1})$.

The operational objective is to produce, for each relevant change event,
an \textbf{Action Ticket}

\[
A = (\text{summary}, \text{risk\_score},
\text{recommended\_actions}, \text{owner},
\text{urgency}, \text{evidence\_citations})
\]

that is grounded in retrieved context and auditable by human analysts.

The gap in current approaches is threefold. First, \textbf{detection
latency}: manual monitoring introduces hours-to-days delay between a
platform change and organizational awareness. Second, \textbf{triage
noise}: without semantic filtering, teams are overwhelmed by cosmetic
changes (formatting, navigation updates) that are operationally
irrelevant. Third, \textbf{recommendation drift}: without grounding in
retrieved evidence, LLM-generated recommendations may hallucinate
actions that are inconsistent with actual platform behavior.

\subsection{ECI Framework}\label{eci-framework}

Ecosystem Change Intelligence (ECI) addresses this gap through three
formal stages:

\textbf{Stage 1: Delta Detection.} For each source $s_i$, the system
periodically fetches a new snapshot $s_i(t_k)$, computes a structured diff
against the previous snapshot $s_i(t_{k-1})$, and produces a change event
$\Delta_i(t)$ containing the added text, deleted text, change type, and
source metadata. This differs from naive textual diffing in that it
applies HTML cleaning, boilerplate removal, and semantic deduplication
before comparison.

\textbf{Stage 2: Change-Grounded Retrieval.} Only the delta content
(added and deleted spans) is chunked (1,600 characters, 200-character
overlap), embedded using a task-aware model with \texttt{search\_document:}
prefix optimization, and stored in a vector index. Restricting the
embedding pipeline to delta content (rather than embedding full
documents) is a deliberate design choice: it concentrates the index on
operationally relevant signal and avoids the vector dilution that
plagues full-document RAG over long, repetitive ecosystem corpora. In
parallel, a knowledge graph (KG) layer indexes entities extracted from
each change event (CVEs, APIs, permissions, policy clauses) and links
them to their originating change. At query time, the system performs
vector retrieval against the delta index and a 2-hop KG traversal, and
fuses the two ranked lists via Reciprocal Rank Fusion. We refer to this
fused retrieval methodology as \textbf{DeltaRAG}.

\textbf{Stage 3: Evidence-Backed Action Generation.} For changes deemed
relevant by the Sentinel Agent (relevance score $\geq$ threshold), a
Coordinator Agent consumes the DeltaRAG output, identifies cross-source
connections (e.g., a CVE linked to an API deprecation linked to a policy
change), and generates Action Tickets with structured fields: summary,
risk score, recommended actions, priority, owner suggestion, and
evidence citations. All outputs are constrained to a strict JSON schema
to ensure downstream parseability.

\subsection{Assumptions}\label{assumptions}

The ECI framework operates under the following assumptions: (1) Public
ecosystem sources reflect genuine signal changes relevant to operational
risk; (2) Textual diffs between cleaned snapshots capture meaningful
platform updates; (3) Semantic similarity in the embedding space
approximates contextual relevance for risk assessment; (4) Composite
risk scoring (combining relevance, severity, and affected signal count)
can approximate fraud impact; (5) Lightweight instruction-tuned LLMs
(7--8B parameters) can produce reliable structured JSON outputs with
appropriate prompting.

\subsection{Formal Framework}\label{formal-framework}

We now formalize the core ECI operations. Let
$\mathcal{S} = \{s_1, s_2, \ldots, s_n\}$ be the set of monitored
sources. Each source $s_i$ produces a time-ordered snapshot sequence:
\begin{equation}
  \mathcal{T}_i = \langle s_i(t_1),\, s_i(t_2),\, \ldots \rangle, \quad t_1 < t_2 < \cdots
\end{equation}

\noindent\textbf{Delta Detection.} A \emph{change event} is the
structured difference between consecutive snapshots:
\begin{equation}
  \Delta_i(t_k) = \mathrm{Diff}\!\left(s_i(t_k),\, s_i(t_{k-1})\right)
                = \bigl\{(\mathrm{kind}_j,\, \mathrm{text}_j)\bigr\}_{j=1}^{m}
\end{equation}
where $\mathrm{kind}_j \in \{\mathrm{added}, \mathrm{deleted}, \mathrm{context}\}$ and cosmetic tokens are removed before comparison.

\noindent\textbf{Chunking.} The delta is segmented into overlapping spans:
\begin{equation}
  \mathcal{C}\!\left(\Delta_i(t_k)\right) = \{c_1, c_2, \ldots, c_p\}
\end{equation}
where $|c_j| \leq L_{\max} = 1600$ characters and consecutive chunks share $L_{\mathrm{ovlp}} = 200$ characters of overlap.

\noindent\textbf{Task-Aware Embedding.} Each chunk is embedded with a retrieval-optimized prefix:
\begin{equation}
  \mathbf{e}_j = \phi\!\left(\text{``search\_document: ''} \mathbin\| c_j\right) \in \mathbb{R}^{768}
\end{equation}
where $\phi$ is the Nomic Embed v1.5 encoder and $\mathbin\|$ denotes string concatenation. The chunking and embedding pipeline (Eqs. 3 and 4) is applied only to delta content; the resulting vector index $\mathcal{D}$ therefore contains only chunks from $\mathrm{kind}_j \in \{\mathrm{added}, \mathrm{deleted}\}$.

\noindent\textbf{Vector Retrieval.} Both the Standard RAG baseline and DeltaRAG query the same delta-only vector index:
\begin{equation}
  R_{\mathrm{vec}}(q,k) =
    \operatorname*{top\text{-}k}_{j \in \mathcal{D}}\,
    \cos\!\Bigl(\phi_q(q),\, \mathbf{e}_j\Bigr)
\end{equation}
where $\phi_q(q) = \phi(\text{``search\_query: ''} \mathbin\| q)$. The Standard RAG variant evaluated in Section~\ref{ablation-study} corresponds exactly to $R_{\mathrm{vec}}$ without any further augmentation.

\noindent\textbf{Knowledge Graph.} The KG layer maintains:
\begin{equation}
  \mathcal{G} = (V, E), \quad
  V = \mathcal{V}_{\mathrm{ent}} \cup \mathcal{V}_{\mathrm{chg}}, \quad
  E \subseteq V \times V \times R
\end{equation}
where $\mathcal{V}_{\mathrm{ent}}$ are extracted entities (CVE IDs, API levels, permission names, policy clauses), $\mathcal{V}_{\mathrm{chg}}$ are change event nodes, and $R = \{\textit{affects}, \textit{deprecates}, \textit{co-occurs}, \textit{references}, \textit{supersedes}\}$.

\noindent\textbf{2-Hop Neighbourhood.} The KG expands retrieval context via bounded traversal:
\begin{equation}
  N_2(v) = \bigl\{u \in V : d_{\mathcal{G}}(v, u) \leq 2\bigr\}
\end{equation}
where $d_{\mathcal{G}}$ is the unweighted shortest-path distance in $\mathcal{G}$.

\noindent\textbf{KG Context Expansion.}
\begin{equation}
  \mathrm{KGCtx}(q) =
    \bigcup_{v \in \mathcal{E}(q)} \bigl\{c_j : c_j \text{ is linked to } N_2(v)\bigr\}
\end{equation}
where $\mathcal{E}(q)$ extracts named entities from query $q$.

\noindent\textbf{Reciprocal Rank Fusion.} The two ranked lists are merged via RRF:
\begin{equation}
  \mathrm{RRF}(d) = \sum_{r \in \mathcal{R}} \frac{1}{k_{\mathrm{rrf}} + \mathrm{rank}_r(d)}, \quad k_{\mathrm{rrf}} = 60
\end{equation}
where $\mathcal{R}$ is the set of input ranked lists.

\noindent\textbf{DeltaRAG (Full System).} DeltaRAG fuses the vector retrieval with the KG context expansion via RRF:
\begin{equation}
  \mathrm{DeltaRAG}(q,k) =
    \operatorname*{top\text{-}k}_{d}\,\mathrm{RRF}(d),
    \quad d \in R_{\mathrm{vec}}(q) \cup \mathrm{KGCtx}(q)
\end{equation}
The contribution of DeltaRAG over Standard RAG is thus the addition of the KG component and the RRF fusion step; the vector half is shared.

\noindent\textbf{Composite Risk Score.} The Coordinator Agent evaluates the synthesized context to assign a scalar risk score to each Action Ticket. Rather than a rigid heuristic, the LLM-driven Coordinator approximates a composite function:
\begin{equation}
  \rho(a) \approx \alpha \cdot r_{\mathrm{rel}}(a) + \beta \cdot r_{\mathrm{sev}}(a) + \gamma \cdot r_{\mathrm{scope}}(a)
\end{equation}
where $r_{\mathrm{rel}} \in [0,10]$ is the upstream Sentinel relevance, $r_{\mathrm{sev}} \in [0,10]$ is the extracted vulnerability severity, and $r_{\mathrm{scope}}$ scales with the volume of affected cross-source signals.

\section{System Architecture}\label{system-architecture}

The ECI system is organized as an six-stage pipeline spanning three architectural layers: the Platform Evolution Layer (data acquisition and change detection), the Semantic Retrieval Layer (embedding, storage, and similarity search), and the Agent Intelligence Layer (triage, coordination, and recommendation generation). We describe each component below.

\subsection{Platform Evolution Layer}\label{platform-evolution-layer}

\emph{\textbf{Source Registry}}

Monitored sources are stored in a relational registry table in Supabase. Each record contains the source URL, name, fetch type (HTML or API/JSON), source category (\texttt{security\_bulletin}, \texttt{developer\_docs}, \texttt{cve\_feed}, \texttt{policy\_update}), and active status. The registry uses an idempotent upsert on the URL field, ensuring re-seeding the pipeline does not create duplicate entries.

\emph{\textbf{Snapshot Ingestion}}

Automated scripts periodically fetch raw content from each active source. HTML pages are retrieved via \texttt{requests} and parsed with \texttt{BeautifulSoup4} to strip navigation, headers, footers, and scripts, retaining only the structural body content. API/JSON sources (e.g., the CISA KEV feed) are fetched directly and serialized. Each fetch produces a snapshot record containing the raw text, cleaned text, content hash (SHA-256), and timestamp, enabling full historical lineage.

\emph{\textbf{Change Detection}}

For each new snapshot, the system retrieves the most recent prior snapshot for the same source and computes a structured diff. The diffing logic identifies added and deleted blocks. Trivial, cosmetic changes (whitespace normalization, minor formatting shifts) are filtered using a strict lexical distance threshold computed via \texttt{difflib.SequenceMatcher}. Changes that exceed the threshold are recorded as change events in a dedicated \texttt{changes} table, linked to both the previous and new snapshot IDs.

\subsection{Semantic Retrieval Layer}\label{semantic-retrieval-layer-deltarag}

\emph{\textbf{Chunking Strategy}}

Change content is chunked into 1,600-character segments with 200-character overlaps. This 1600/200 rule balances granularity (ensuring each chunk is semantically focused) against context preservation (preventing sentence breaks). Before embedding, each chunk is prefixed with \texttt{search\_document:} to activate Nomic v1.5's task-aware encoding path, optimizing the embedding strictly for retrieval rather than classification.

\emph{\textbf{Embedding Model: Nomic v1.5}}

We selected Nomic Embed v1.5 based on strict architectural constraints for processing security data. It offers three critical advantages over standard models: (1) 768-dimensional embeddings that accurately capture nuanced distinctions between vulnerability classes; (2) an 8,192-token maximum input length, accommodating long JSON feeds and policy documents without truncating context; and (3) local execution capability, ensuring zero-day vulnerability data never leaves the organization's controlled infrastructure.

\emph{\textbf{Vector Storage: Supabase + pgvector}}

Embedded chunks are stored in a Supabase PostgreSQL database utilizing the \texttt{pgvector} extension. The \texttt{embeddings} table stores: chunk ID, source ID, source category, change ID, chunk kind (addition, deletion), chunk text, and a \texttt{vector(768)} embedding column. An HNSW (Hierarchical Navigable Small World) index enables sub-10ms cosine similarity search. Using PostgreSQL-native storage rather than an isolated vector database (e.g., Pinecone) enables seamless SQL joins between vector representations, relational change events, and downstream agent recommendations.

\subsection{Agent Intelligence Layer}\label{agent-intelligence-layer}

\emph{\textbf{Sentinel Triage Agent}}

The Sentinel Agent serves as the first-pass filter. For each new change event, it receives the structured change description and performs semantic comprehension, identifying what the change means in the context of ecosystem fraud and security operations. It assigns a relevance score (threshold $\geq5$) and classifies the domain (e.g., \texttt{vulnerability\_exposure}). Most importantly, it distinguishes signal from noise, acting as a gatekeeper to prevent the Coordinator Agent from processing irrelevant documentation updates. The Sentinel outputs strict, validated JSON.

\emph{\textbf{Coordinator Agent (KG-Augmented Retrieval)}}

Changes that pass the Sentinel's threshold escalate to the Coordinator
Agent for deep analysis. The Coordinator performs the full
\textbf{DeltaRAG} retrieval (Eq.~10): vector search over the
delta-chunk index is executed in parallel with a 2-hop traversal of
the knowledge graph (KG), and the two ranked lists are fused via
Reciprocal Rank Fusion. The KG construction and traversal details are
described in Section~\ref{graph-rag-for-cross-source-intelligence}.
From the fused output, the Coordinator generates structured insights
explaining why the change matters, aggregates related high-risk events
(e.g., linking a CISA KEV alert to an OEM security bulletin), and
produces Action Tickets with priority levels, suggested owners, and
explicit evidence citations pointing directly to the retrieved chunks.

\emph{\textbf{Pipeline Orchestration \& Logging}}

Rather than relying on heavy multi-agent orchestration frameworks, the system employs an asynchronous backend worker (\texttt{run\_pipeline.py}) that polls a \texttt{pipeline\_jobs} queue. As the worker executes the sequential stages (Ingest $\rightarrow$ Diff $\rightarrow$ Vectorize $\rightarrow$ KG Update $\rightarrow$ Sentinel $\rightarrow$ Coordinator), it streams \texttt{stdout} natively into a relational \texttt{pipeline\_logs} table. This architecture guarantees full traceability of every retrieval step and LLM invocation, driving the real-time execution monitor on the Next.js dashboard.


\section{Knowledge Graph for Cross-Source Intelligence}\label{graph-rag-for-cross-source-intelligence}

\subsection{Knowledge Graph Construction}\label{knowledge-graph-construction}

The knowledge graph (KG) layer constructs a directed graph from change
events and their extracted metadata. Nodes represent distinct entities
(CVE identifiers, Android API levels, permission names, SDK versions,
policy clause identifiers, components) and change events. Edges encode
the structural relationship set $R$ defined in Eq.~6: $\{\textit{affects},
\textit{deprecates}, \textit{co-occurs}, \textit{references},
\textit{supersedes}\}$. To ensure deterministic execution and minimize
latency, entity extraction is performed using a pure regex-based
extractor (\texttt{rag/entity\_extractor.py}) over strict pattern
matching and predefined whitelists. This design trades extraction
coverage for two operational guarantees: zero hallucination at the
graph-construction step, and sub-millisecond entity recognition per
change event.

\subsection{KG-Augmented Retrieval via RRF}\label{graph-augmented-retrieval}

At query time, the system executes a parallel retrieval strategy to
capture both semantic similarity and structural dependencies. Given a
query (e.g., a new bulletin diff):
\begin{enumerate}
    \item \textbf{Vector Search}: Retrieves the top-$K$ semantically similar chunks across all other sources from the delta-chunk index (Eq.~5), using Nomic Embed v1.5.
    \item \textbf{KG Traversal}: Extracts entities from the query, traverses the KG up to 2 hops (Eq.~7) to find connected change events, and gathers their corresponding chunks (Eq.~8).
\end{enumerate}
Rather than naively appending these lists or aggressively boosting graph
distances (which would displace highly-relevant semantic matches), the
system merges the two sets using Reciprocal Rank Fusion (RRF, Eq.~9)
with $k_{\mathrm{rrf}} = 60$. This ensures that a chunk ranked highly by
Vector Search retains its score even if the KG fails to find an entity
path to it, providing robust fallback behavior. The fused output is the
\textbf{DeltaRAG} retrieval defined in Eq.~10.

\subsection{Entity-Weighted Scoring \& Cross-Category Bonus}\label{adaptive-retrieval-strategy-learning}

A naive graph traversal in a highly interconnected corpus quickly
succumbs to noise: connector entities such as \texttt{android\_14} are
shared across hundreds of unrelated changes and would dominate
retrieval if treated identically to rigid identifiers. To suppress this,
the traversal engine applies \textbf{entity-type weighting}: nodes are
scored in proportion to their discriminative power. Unique vulnerability
identifiers (\texttt{cve}, including SVE-prefixed entries) carry a
$3.0\times$ weight; domain-specific terms (\texttt{policy\_clause},
\texttt{permission}) carry $2.0\times$; and generic structural nodes
(\texttt{component}, \texttt{api\_level}, \texttt{kernel\_version},
\texttt{sdk\_version}) carry $1.0\times$. A change's score aggregates
the weighted entities it shares with the query, divided by the minimum
hop distance:
\begin{equation}
  \mathrm{score}(c) = \frac{\sum_{e \in \mathrm{shared}(c, q)} w_e}{1 + h_{\min}(c, q)}
\end{equation}
where $w_e$ is the entity-type weight of $e$ and $h_{\min}(c, q)$ is the
minimum hop distance from any query entity to change $c$ in $\mathcal{G}$.

Finally, the system applies a $1.5\times$ \textbf{cross-category bonus}
to changes whose source category differs from the query's source
category, rewarding the graph for surfacing connections that bridge
distinct source families (e.g., linking a Security Bulletin to an OEM
Patch). This bias is what enables the KG layer to recover the
cross-source linkages that pure vector similarity misses.


\section{Implementation}\label{implementation}

\subsection{Technology Stack}\label{technology-stack}

The system is implemented in Python 3.11 with the following core components. Data acquisition uses \texttt{requests} and \texttt{BeautifulSoup4} for HTML parsing. The Supabase Python client manages all PostgreSQL database interactions (source registry, snapshots, changes, agent events, vector chunks, recommendations). Embedding generation uses the Nomic Embed v1.5 model running locally via the \texttt{sentence-transformers} package. LLM inference currently utilizes Groq's API for ultra-low-latency access to \texttt{llama-3.1-8b-instant}. The presentation layer is a custom Next.js application utilizing React, Tailwind CSS, and Framer Motion, interfacing directly with the Supabase backend.

\subsection{Database Schema}\label{database-schema}

The relational schema in Supabase comprises core tables tracking the lifecycle of an intelligence event: \texttt{sources} (registry of monitored URLs), \texttt{snapshots} (versioned page content with content hashes), \texttt{changes} (structured diff records linking previous and new snapshots), \texttt{embeddings} (pgvector chunks with lineage metadata), \texttt{agent\_events} (per-change-event triage analysis records), and \texttt{recommendations} (final Action Tickets). Two operational tables, \texttt{pipeline\_jobs} (worker queue) and
\texttt{pipeline\_logs} (streamed worker stdout), support the asynchronous
orchestration described in Section~\ref{agent-intelligence-layer}. Foreign key relationships enforce referential integrity across the pipeline: changes reference snapshots, embeddings reference changes and sources, and recommendations reference specific change events. The \texttt{knowledge\_graph\_data} table acts as a JSONB persistence layer for the graph.

\subsection{LLM Selection}\label{llm-evaluation-and-selection}

The Sentinel and Coordinator roles require strict JSON stability and sub-second inference times. We selected \texttt{llama-3.1-8b-instant} served via the Groq API as the primary inference engine. This model achieves the optimal balance of instruction-following capabilities (required for the Coordinator's complex, nested evidence citations) and ultra-low latency. All outputs are subjected to strict parse validation: outputs must pass \texttt{json.loads()} and conform to predefined Pydantic schemas without manual post-processing.

\subsection{Prompt Engineering for Structured Output}\label{prompt-engineering-for-structured-output}

Agent prompts are carefully engineered to maximize JSON stability and minimize hallucination. The Sentinel prompt includes: (1) a system-level role definition constraining the agent to a security analyst persona; (2) the complete JSON schema with field descriptions and value constraints; (3) explicit instructions to output only the JSON object with no preamble or explanation; and (4) few-shot examples demonstrating correct output for both relevant and irrelevant changes. The Coordinator prompt dynamically integrates retrieved context chunks (formatted as numbered evidence blocks) and enforces explicit citation instructions, requiring \texttt{evidence\_ids} for each synthesized recommendation.

\section{Evaluation}\label{preliminary-evaluation}

\subsection{Embedding Model Selection}\label{retrieval-quality}

The ECI pipeline processes massive, highly technical documents including JSON vulnerability feeds and multi-page OEM security bulletins. Standard embedding models often struggle with these formats due to strict token limits, leading to truncation of critical security context.

We evaluated candidate embedding models against three non-negotiable architectural constraints for the ECI pipeline:
\begin{enumerate}
    \item \textbf{High Context Window}: Must support $>4,000$ tokens to embed large policy updates and CISA KEV arrays without aggressive chunking.
    \item \textbf{Data Privacy}: Must run locally to ensure unreleased, proprietary OEM vulnerability data is not leaked to third-party cloud APIs.
    \item \textbf{Asymmetric Retrieval}: Must support task-aware representations to match short queries (diffs) against long documents (bulletins).
\end{enumerate}

Based on these constraints, we selected \textbf{Nomic-Embed-Text-v1.5}.

{\def\LTcaptype{none}
\begin{longtable}[]{@{}
  >{\centering\arraybackslash}p{(\linewidth - 8\tabcolsep) * \real{0.2619}}
  >{\centering\arraybackslash}p{(\linewidth - 8\tabcolsep) * \real{0.1529}}
  >{\centering\arraybackslash}p{(\linewidth - 8\tabcolsep) * \real{0.1747}}
  >{\centering\arraybackslash}p{(\linewidth - 8\tabcolsep) * \real{0.2577}}@{}}
\toprule\noalign{}
\textbf{Model} & \textbf{Privacy} & \textbf{Max Tokens} & \textbf{Decision} \\
\midrule\noalign{}
\endhead
\bottomrule\noalign{}
\endlastfoot
Nomic v1.5 & Local & 8,192 & Selected \\
BAAI BGE-Large & Local & 512 & Ruled out (Truncation) \\
OpenAI / Vertex AI & Cloud & 8,192+ & Ruled out (Privacy) \\
\end{longtable}
}

\emph{Table 1: Embedding model selection matrix based on ECI architectural constraints.}

Nomic v1.5's primary advantage stems from its task-aware prefix mechanism (\texttt{search\_document:} / \texttt{search\_query:}), 768-dimensional output, and 8,192-token context window that successfully avoids the truncation which degrades models like BGE on long policy documents.

\subsection{Ablation Study: Retrieval Variants}
\label{ablation-study}

To isolate the contribution of knowledge graph (KG) augmentation within DeltaRAG, we conducted a systematic ablation across three retrieval variants evaluated on a benchmark of 100 deterministic queries spanning ten query types
(10 queries each). Each query is anchored to one of the ten seeded
sources, with deterministic cross-source ground truth derived from
shared-entity co-occurrence in the synthetic snapshot data. The ten
query types probe distinct retrieval regimes:

\begin{description}
  \item[\textit{full\_diff} (A)] The raw \texttt{diff\_text} of a change
    used verbatim as the query, this mirrors the pipeline's own
    internal trigger and represents the upper bound on context richness.
  \item[\textit{identifier\_focused} (B)] CVE / SVE / API-identifier
    extracts (e.g.\ \texttt{CVE-2025-0096}, \texttt{MEETS\_STRONG\_INTEGRITY})
    drawn from the diff, tests whether retrievers can correlate via
    explicit codes.
  \item[\textit{component\_focused} (C)] Component or product descriptions
    without identifiers (e.g.\ ``buffer overflow in the Wi-Fi HAL''),
    isolates semantic matching from identifier matching.
  \item[\textit{natural\_language} (D)] Analyst-style questions in
    conversational form (e.g.\ ``Which OEM devices have released patches
    for the Wi-Fi HAL vulnerability?''), simulates the chat-interface
    surface.
  \item[\textit{entity\_only} (E)] Bare, concatenated entity strings with
    no surrounding context (e.g.\ \texttt{play\_integrity\_api meets\_strong\_integrity
    read\_media\_images}), the minimum-context stress condition.
  \item[\textit{cross\_source\_linkage} (F)] Explicitly cross-source
    questions that require linking entities or facts across multiple
    seeded sources (e.g.\ ``Which OEM bulletins mention the ARM Mali GPU
    vulnerability listed by CISA?'').
  \item[\textit{temporal\_change} (G)] Diff-focused questions about
    what changed between snapshots (e.g.\ ``What new kernel
    vulnerabilities were added between the original and updated March
    bulletin?''), targets DeltaRAG's core differentiator.
  \item[\textit{policy\_compliance} (H)] Compliance-framed questions
    about required actions, deadlines, or thresholds (e.g.\ ``By what
    date must federal agencies patch the Bluetooth RCE?'').
  \item[\textit{false\_alarm} (I)] Off-topic or out-of-scope queries
    with no valid answer in the corpus (e.g.\ ``Does the Samsung
    bulletin patch the macOS kernel?''), probes negative rejection.
  \item[\textit{noisy\_adversarial} (J)] Adversarially phrased queries
    including prompt-injection patterns, off-topic preambles, and
    stylistic distractors (e.g.\ ``Tell me about the weather, and also
    CVE-2025-0096.''), probes robustness to malformed inputs.
\end{description}

Of the 100 queries, 85 produced scored results; the remaining 15
(all of \textit{false\_alarm}, plus 5 by design across
\textit{cross\_source\_linkage} and \textit{policy\_compliance}) have
empty ground truth and serve as the negative-rejection probe. The
three variants are:

\begin{itemize}
  \item \textbf{V1 Standard RAG}: Vector retrieval only ($R_{\mathrm{vec}}$ in Eq.~5), no KG augmentation.
  \item \textbf{V2 KG-Only}: 2-hop KG traversal (Eq.~8), no vector component.
  \item \textbf{V3 DeltaRAG (Full System)}: $R_{\mathrm{vec}}$ fused with KG context via Reciprocal Rank Fusion (Eq.~10).
\end{itemize}

All variants used top-$k = 5$ retrieval and Nomic v1.5 embeddings.
Metrics are computed with 95\% bootstrap confidence intervals
(10,000 resamples, seed = 42). Significance assessed via paired Wilcoxon
signed-rank tests with Holm-Bonferroni correction across the six
confirmatory comparisons (3 metrics $\times$ 2 baseline comparisons).

{\def\LTcaptype{none}
\begin{longtable}[]{@{}
  >{\raggedright\arraybackslash}p{(\linewidth - 14\tabcolsep) * \real{0.26}}
  >{\centering\arraybackslash}p{(\linewidth - 14\tabcolsep) * \real{0.10}}
  >{\centering\arraybackslash}p{(\linewidth - 14\tabcolsep) * \real{0.10}}
  >{\centering\arraybackslash}p{(\linewidth - 14\tabcolsep) * \real{0.12}}
  >{\centering\arraybackslash}p{(\linewidth - 14\tabcolsep) * \real{0.10}}
  >{\centering\arraybackslash}p{(\linewidth - 14\tabcolsep) * \real{0.10}}
  >{\centering\arraybackslash}p{(\linewidth - 14\tabcolsep) * \real{0.10}}
  >{\centering\arraybackslash}p{(\linewidth - 14\tabcolsep) * \real{0.10}}@{}}
\toprule\noalign{}
\textbf{Variant} & \textbf{P@1} & \textbf{P@5} & \textbf{R@5} & \textbf{nDCG@5} & \textbf{MRR} & \textbf{MAP} & \textbf{FA-Rej} \\
\midrule\noalign{}
\endhead
\bottomrule\noalign{}
\endlastfoot
V1 Standard RAG                  & 0.847          & \textbf{0.504} & \textbf{0.931} & 0.877          & 0.913          & \textbf{0.822} & 1.000 \\
V2 KG-Only                & 0.635          & 0.341          & 0.571          & 0.587          & 0.639          & 0.563          & 1.000 \\
V3 DeltaRAG (Full)   & \textbf{0.882} & 0.478          & 0.899          & \textbf{0.873$^*$} & \textbf{0.930} & 0.821          & 1.000 \\
\end{longtable}
}

\emph{Table 2: Ablation results across 85 scored benchmark queries, top-$k$\,=\,5. Bold marks the best variant per metric. $^*$V3 nDCG@5 is statistically indistinguishable from V1 ($p = 0.43$); V3 wins on P@1 and MRR.}

\textbf{Headline numbers (Holm-corrected, $\alpha = 0.05$):}
\begin{itemize}
  \item \textbf{V3 vs.\ V2 (KG-Only)}: V3 strictly dominates on every metric, $\Delta$R@5 = +0.328 [+0.235, +0.423], $\Delta$nDCG@5 = +0.286 [+0.204, +0.373], both $p < 0.001$.
  \item \textbf{V3 vs.\ V1 (Standard RAG)}: V3 wins on top-rank quality (\textbf{P@1} = 0.882 vs.\ 0.847; \textbf{MRR} = 0.930 vs.\ 0.913) and ties on nDCG@5 ($p = 0.43$). V1 retains a small edge on raw coverage, R@5 ($\Delta = -0.032$, $p = 0.014$) and P@5 ($\Delta = -0.026$, $p = 0.008$).
\end{itemize}

The per-type breakdown reveals where the KG component pays off: V3
strictly outperforms V1 on \textit{policy\_compliance} (R@5 0.875 vs.\ 0.833;
nDCG@5 0.685 vs.\ 0.653), edges ahead on \textit{temporal\_change} nDCG@5
(0.772 vs.\ 0.750) at identical recall, and beats V1 on
\textit{noisy\_adversarial} nDCG@5 (0.766 vs.\ 0.749), early evidence that
the KG component dampens rank instability under adversarial phrasing.

The most operationally relevant finding is the P@1 and MRR uplift: V3
places a relevant result at rank 1 on 88.2\% of scored queries
(vs.\ 84.7\% for V1), and its reciprocal rank is 0.930 (vs.\ 0.913). For a
triage system whose downstream consumer is an analyst clicking the top
result, this top-of-list precision matters more than depth-5 recall. The
dense embedding model already correlates explicit cross-source identifiers
(e.g.\ \texttt{CVE-2025-0096}) effectively enough that pure-vector retrieval
is hard to beat on raw coverage, but the KG-augmented fusion provides
the structural signal needed to push the correct answer to the top of the
list. The standalone KG-Only variant (V2) remains constrained by a
regex-based entity extraction bottleneck on free-text vulnerability
descriptions; addressing this via LLM-driven entity extraction is a
critical focus for future iterations.

\section{Discussion}
\label{discussionn}

\subsection{DeltaRAG vs. Standard RAG}\label{deltarag-vs.-standard-rag}

DeltaRAG's design rests on two distinct decisions, one at the pipeline
level and one at the retrieval level. The pipeline-level decision is to
embed only delta content (additions and deletions) rather than full
documents. This concentrates the vector index on operationally
relevant signal and is shared by the Standard RAG baseline (V1) in our
ablation, both variants query the same delta-only chunk store.
The retrieval-level decision is the KG augmentation: at query time,
vector retrieval is fused with a 2-hop knowledge graph traversal via
Reciprocal Rank Fusion (Eq.~10). This is what distinguishes V3
DeltaRAG from V1 Standard RAG.

Empirically (Section~\ref{ablation-study}), this yields a small but
consistent uplift in top-rank quality (P@1\,$+0.035$, MRR\,$+0.017$)
over V1 while remaining statistically tied on nDCG@5, a favorable
trade for triage workflows where an analyst acts on the first surfaced
result. The KG layer's contribution is most visible on
\textit{policy\_compliance}, \textit{temporal\_change}, and
\textit{noisy\_adversarial} queries, where structural linkage matters
more than dense semantic similarity.

\subsection{Knowledge Graph for Cascading Risk Detection}\label{graph-rag-for-cascading-risk-detection}

The knowledge graph (KG) layer addresses a fundamental limitation of
vector-only retrieval: the inability to surface connections that are
semantically distant but structurally linked. For example, a CVE
affecting the Android kernel (from the CISA KEV feed) and a Google
Play Developer Policy update regarding background data access may not
share similar vocabulary. However, their operational intersection is
critical. The KG captures this by linking both changes to explicit
shared entities, such as a specific \texttt{API\_level} or a shared
system component like \texttt{ActivityManagerService}. This strict
structural linkage enables the Coordinator Agent to surface the
connection and recommend coordinated action without relying on vector
proximity.

\subsection{Limitations and Threats to Validity}\label{limitations-and-threats-to-validity}

Several limitations should be noted. First, while our ablation study
utilized 100 deterministic queries (85 scored) with paired non-parametric
significance testing, Holm-Bonferroni multiple-comparison correction, and
10,000-resample bootstrap confidence intervals, the evaluation was conducted
against a synthetically seeded snapshot database. True differentiation
between hybrid retrieval mechanisms requires independent queries annotated
against live, stream-processed platform diffs. Second, per-type sample size
is modest ($n = 6$--$10$), which limits statistical power for the
exploratory per-type comparisons; effect sizes and bootstrap CIs are
therefore reported alongside raw $p$-values. Third, the Knowledge Graph
currently relies entirely on regex-based entity extraction. While this
guarantees high execution speed and zero hallucination, it introduces a
critical bottleneck: it misses implicit, unstructured relationships that
do not conform to predefined string patterns, and this directly accounts
for the underperformance of the standalone KG-Only variant (V2)
reported in Section 7.2. Fourth, lightweight instruction-tuned LLMs (e.g.,
Llama 3.1 8B) trade deep reasoning capabilities for latency and JSON parse
stability; complex, multi-stage risk scenarios ultimately require
human-in-the-loop validation.


\subsection{Operational Implications}\label{operational-implications}

The ECI framework is designed to reduce three operational costs: (1) time-to-awareness, by detecting platform changes within minutes rather than days; (2) triage burden, by automatically filtering irrelevant cosmetic changes and surfacing only risk-relevant updates via the Sentinel Agent; and (3) decision latency, by generating evidence-backed Action Tickets that can be directly routed to the appropriate operational queue with assigned priorities. For fraud operations teams at financial institutions, this translates to faster signal recalibration, reduced fraud exposure windows, and drastically decreased manual monitoring effort.

\section{Conclusion and Future
Work}\label{conclusion-and-future-work}

We have presented Ecosystem Change Intelligence (ECI), a framework for
converting ecosystem monitoring from a passive documentation-tracking
activity into an active, change-centric decision pipeline. Our primary
contribution, \textbf{DeltaRAG} (Full System), grounds
retrieval on change-local context windows (Eq.~5) and fuses them with
2-hop knowledge graph traversal via RRF, achieving
\textbf{P@1\,=\,0.882}, \textbf{MRR\,=\,0.930}, \textbf{nDCG@5\,=\,0.873},
and \textbf{100\% false-alarm rejection} across 100 benchmark queries
(85 scored). On the metrics most relevant to a triage workflow,
top-rank quality, the Full System outperforms Standard RAG (P@1
$+0.035$, MRR $+0.017$) while remaining statistically tied on nDCG@5
($p = 0.43$); it strictly dominates the KG-Only baseline on every
metric ($\Delta$R@5 $= +0.328$, $p < 0.001$). The per-type breakdown
identifies the regimes where KG augmentation pays off most clearly, \textit{policy\_compliance}, \textit{temporal\_change}, and \textit{noisy\_adversarial} queries. Complementary contributions include a formal mathematical framework (Eqs.~1--11), a KG layer for cross-source evidence linking, a two-agent architecture (Sentinel + Coordinator) producing auditable Action Tickets with composite risk scoring, and an ablation study across three retrieval variants validating the design choices. The system monitors 10 live
Android ecosystem sources and surfaces detected changes to fraud operations teams as evidence-backed, priority-ranked recommendations.

Future work will pursue four directions. First, \textbf{benchmark
quality}: replacing the current synthetic benchmark with independent
queries annotated against live scraped diffs, and broadening the
robustness probe to include stale-information retrieval alongside the
existing noisy-adversarial and false-alarm categories. Second,
\textbf{retrieval improvement}: addressing the regex-based entity
extraction bottleneck (the dominant cause of V2 KG-Only
underperformance) via LLM-driven extraction, and implementing BM25
hybrid search and cross-encoder reranking to lift the current
P@5\,=\,0.478. Third, \textbf{validation}: conducting analyst
walkthroughs to evaluate Action Ticket quality in real fraud operations
workflows, and iterating on risk scoring weights ($\alpha, \beta,
\gamma$ in Eq.~11) based on operational feedback. Fourth,
\textbf{infrastructure}: migrating LLM inference to self-hosted GCP
with NVIDIA NIM/vLLM for production-grade performance and data privacy,
and hardening structured outputs through prompt optimization and schema
validation layers.

\section{References}\label{references}

{[}1{]} Aghaei, E., et al. (2023). SecureBERT: A Domain-Specific
Language Model for Cybersecurity. In Proceedings of ACM CCS Workshop on
Artificial Intelligence and Security.

{[}2{]} Edge, D., et al. (2024). From Local to Global: A Graph RAG
Approach to Query-Focused Summarization. arXiv:2404.16130.

{[}3{]} Ferrag, M.A., et al. (2024). Generative AI and Large Language
Models for Cyber Security: All Insights You Need. arXiv:2405.12750.

{[}4{]} He, Y., et al. (2024). G-Retriever: Retrieval-Augmented
Generation for Textual Graph Understanding and Question Answering. In
Proceedings of NeurIPS.

{[}5{]} Izacard, G., \& Grave, E. (2021). Leveraging Passage Retrieval
with Generative Models for Open Domain Question Answering. In
Proceedings of EACL.

{[}6{]} Lewis, P., et al. (2020). Retrieval-Augmented Generation for
Knowledge-Intensive NLP Tasks. In Proceedings of NeurIPS.

{[}7{]} Li, G., et al. (2023). CAMEL: Communicative Agents for ``Mind''
Exploration of Large Language Model Society. In Proceedings of NeurIPS.

{[}8{]} Li, Z., et al. (2022). Constructing Knowledge Graphs for Cyber
Threat Intelligence. In Proceedings of IEEE S\&P Workshop.

{[}9{]} Li, Z., et al. (2024). Automated Policy Change Detection and
Tracking: A Survey. ACM Computing Surveys.

{[}10{]} Mittal, S., et al. (2016). CyberTwitter: Using Twitter to
Generate Alerts for Cybersecurity Threats and Vulnerabilities. In
Proceedings of IEEE/ACM ASONAM.

{[}11{]} Nussbaum, Z., et al. (2024). Nomic Embed: Training a
Reproducible Long Context Text Embedder. arXiv:2402.01613.

{[}12{]} Patel, A., et al. (2023). Web Content Change Detection:
Methods, Challenges, and Applications. Information Processing \&
Management, 60(3).

{[}13{]} Reimers, N., \& Gurevych, I. (2019). Sentence-BERT: Sentence
Embeddings using Siamese BERT-Networks. In Proceedings of EMNLP.

{[}14{]} Savelka, J., et al. (2023). Explaining Legal Concepts with
Augmented Large Language Models. In Proceedings of ICAIL.

{[}15{]} Shi, W., et al. (2023). REPLUG: Retrieval-Augmented Black-Box
Language Models. In Proceedings of NAACL.

{[}16{]} Wu, Q., et al. (2023). AutoGen: Enabling Next-Gen LLM
Applications via Multi-Agent Conversation. arXiv:2308.08155.

{[}17{]} Xiao, S., et al. (2024). C-Pack: Packaged Resources to Advance
General Chinese Embedding. In Proceedings of ACL.

{[}18{]} Xiong, G., et al. (2024). Benchmarking Retrieval-Augmented
Generation for Medicine. In Findings of ACL.

{[}19{]} Yao, S., et al. (2023). ReAct: Synergizing Reasoning and Acting
in Language Models. In Proceedings of ICLR.

{[}20{]} Yasunaga, M., et al. (2021). QA-GNN: Reasoning with Language
Models and Knowledge Graphs for Question Answering. In Proceedings of
NAACL.

\end{document}
